\begin{table}[t]
\centering
\caption{Notation used throughout the paper.}
\label{tab:notation}
\begin{tabular}{l p{0.72\columnwidth}}
\toprule
Symbol & Meaning \\
\midrule
$v(t)$ & Membrane voltage (normalized) \\
$r(t)\in[0,1]$ & Continuous reset gate (fast activation-like proxy) \\
$J$ (= $I_{\mathrm{in}}$) & Net input current / drive (normalized) \\
$\tau,\ \tau_r$ & Membrane and gate time constants ($\tau_r\ll\tau$) \\
$v_{\mathrm{rest}},v_{\mathrm{reset}},v_{\mathrm{thresh}}$ & Rest, reset, threshold voltages \\
$\sigma(\cdot)$ & Logistic gate detector, steepness $\alpha$ \\
$\gamma,\ \theta_{\mathrm{eff}}$ & Hysteresis and effective threshold $\theta_{\mathrm{eff}}=v_{\mathrm{thresh}}-\gamma r$ \\
$G_b$ & Fixed reset gain (baseline HyperLIF) \\
$G_{\mathrm{base}},\ \beta$ & Baseline gain and adaptive scaling (Adaptive HyperLIF) \\
$G(r,J)$ & Reset conductance scale; $G_b r^2$ or $(G_{\mathrm{base}}+\beta|J|)r^2$ \\
$G_{\mathrm{dyn}}$ & Adaptive conductance scale $G_{\mathrm{base}}+\beta|J|$ \\
$\lambda$ & Membrane relaxation rate $1/\tau + G$ \\
$J_{\mathrm{DB}}$ & Frozen-open-gate DB branch estimate \\
$\Delta t$ & Simulation time step \\
\bottomrule
\end{tabular}
\end{table}
